% Document généré automatiquement: explication_sans_code.tex
\documentclass[a4paper,11pt]{article}
\usepackage[utf8]{inputenc}
\usepackage[T1]{fontenc}
\usepackage[french]{babel}
\usepackage{graphicx}
\usepackage{amsmath,amssymb}
\usepackage{geometry}
\geometry{margin=2.5cm}
\usepackage{hyperref}

\title{Explication détaillée — Logique floue et ensembles flous (sans code)}
\author{}
\date{\today}

\begin{document}
\maketitle

\section*{Contexte}
Ce document rassemble et explique, sans code, les éléments demandés durant la session : définitions d'ensembles flous, fonctions d'appartenance triangulaires et trapézoïdales, interprétations numériques, exemples concrets (MovieLens) et explications visuelles.

\section{Notations et définitions formelles}
Soit $X$ l'univers (ensemble des valeurs possibles, par exemple l'ensemble des notes). Un ensemble flou $A$ sur $X$ est caractérisé par une fonction d'appartenance
$$\mu_A : X \to [0,1],$$
où $\mu_A(x)$ indique le degré d'appartenance de $x$ à $A$ (0 = pas du tout, 1 = complètement).

Définitions utiles :
\begin{itemize}
  \item Support de $A$ : $\{x\in X\mid \mu_A(x)>0\}$.
  \item Coeur (core) de $A$ : $\{x\in X\mid \mu_A(x)=1\}$.
  \item $\alpha$-cut : $A_\alpha=\{x\in X\mid \mu_A(x)\ge \alpha\}$, pour $\alpha\in(0,1]$.
\end{itemize}

\section{Explication pour un(e) ado de 15 ans}
Imagine que pour chaque valeur (par exemple une note de film) on place un curseur entre 0 et 1 qui indique "à quel point" cette valeur correspond à une idée (par ex. "bon film"). Une fonction d'appartenance donne la position de ce curseur pour chaque valeur. Les fonctions triangulaires et trapézoïdales sont simplement des formes graphiques simples qui définissent ce curseur en fonction de la valeur.

\section{Explication très simple (enfant de 8 ans)}
Pense à une lampe : 0 = éteinte, 1 = très allumée. Si on dit qu'un film est "un peu bon", alors la lampe est à moitié allumée. Un triangle c'est comme une lampe qui atteint son maximum à un seul point (un pic). Un trapèze, c'est une lampe qui reste complètement allumée sur une petite plage (un plateau) puis s'éteint.

\section{Fonctions d'appartenance — triangle et trapèze}
\subsection*{Triangulaire}
Une fonction triangulaire est donnée par trois paramètres $(a,b,c)$, $a<b<c$, et vaut :
$$
\mu(x)=\begin{cases}
0, & x\le a,\\
\dfrac{x-a}{b-a}, & a<x\le b,\\
\dfrac{c-x}{c-b}, & b<x<c,\\
0, & x\ge c.
\end{cases}
$$
Le point $b$ est le sommet où $\mu(b)=1$.

\subsection*{Trapézoïde}
Un trapèze est donné par $(a,b,c,d)$ avec $a\le b\le c\le d$ et vaut :
$$
\mu(x)=\begin{cases}
0, & x\le a,\\
\dfrac{x-a}{b-a}, & a<x\le b,\\
1, & b<x\le c,\\
\dfrac{d-x}{d-c}, & c<x<d,\\
0, & x\ge d.
\end{cases}
$$
Le plateau $[b,c]$ correspond aux valeurs considérées comme pleinement appartenantes.

\section{Exemples concrets et interprétations (MovieLens)}
À partir des données MovieLens analysées, les quantiles observés étaient (exemples) : min = 0.50, q10 = 2.00, q25 = 3.00, q50 = 3.50, q75 = 4.00, q90 = 5.00.

Quelques constructions pratiques utilisées pour illustrer :
\begin{itemize}
  \item Triangles d'exemple : $(2,3,5)$ et $(1,3.5,5)$ (montrent comment le sommet change la sensibilité).
  \item Trapèze d'exemple : $(1,2,4,5)$ (plateau $[2,4]$ où $\mu=1$).
\end{itemize}

Interprétation : superposer ces fonctions sur l'histogramme des notes aide à visualiser quelles plages de notes sont considérées "faibles", "moyennes" ou "fortes" selon la forme choisie. Le trapèze permet de déclarer un intervalle entier comme "completement" appartenant au terme, ce qui correspond à une tolérance humaine.

\section{Questions rapides}
\begin{itemize}
  \item "La somme des degrés vaut-elle 1 ?" Non. Les $\mu$ mesurent indépendamment l'appartenance à chaque terme, ce ne sont pas des probabilités obligatoires.
  \item "Faut-il un triangle ou un trapèze ?"  Trapèze pour tolérance/robustesse, triangle pour précision/focalisation.
  \item "Comment calibrer automatiquement ?" Méthodes courantes : placer paramètres sur des percentiles (ex. q25,q50,q75), ou optimiser par apprentissage pour améliorer les métriques de recommandation.
\end{itemize}

\section*{Conclusion}
Ce document présente, sans code, l'ensemble des explications demandées : formalisme, interprétations intuitives (pour 15 et 8 ans), définitions mathématiques des fonctions triangulaires et trapézoïdales, exemples concrets basés sur MovieLens, et interprétation des figures générées. Si tu souhaites que j'ajoute d'autres sections (défuzzification, moteur Mamdani, exemples chiffrés supplémentaires) dis-moi lesquelles.

\end{document}
